\begin{longtable}{TD}
\caption{Queue-worker lifecycle terms.}
\label{tab:terms-queue-worker-lifecycle}\\

\toprule
Term & Definition \\
\midrule
\endfirsthead

\toprule
Term & Definition \\
\midrule
\endhead

\bottomrule
\endfoot

\termrow{enqueue}{
The lifecycle point at which a task instance is submitted to a queue.
}

\termrow{dequeue}{
The lifecycle point at which a worker takes a task instance from a queue for
processing.
}

\termrow{reserve}{
The lifecycle point at which a worker or worker framework claims a task
instance, often through a lease, visibility timeout, or reservation mechanism.
}

\termrow{start}{
The lifecycle point at which a task attempt begins execution.
}

\termrow{input fetch}{
The stage in which the worker fetches external input, payload data, referenced
objects, or domain state required by the task attempt.
}

\termrow{deserialize}{
The stage in which the worker decodes, parses, or prepares the task payload or
input data for processing.
}

\termrow{processing stage}{
A logical phase inside task execution, such as validation, transformation,
dependency call, write, commit, or cleanup.
}

\termrow{dependency call}{
An interaction with an external service, API, database, storage system, broker,
or remote dependency during task execution.
}

\termrow{retry}{
A repeated execution attempt or repeated operation after failure, timeout,
lease expiration, cancellation, or retryable error.
}

\termrow{backoff}{
A delay before retrying a task attempt or operation. Backoff can contribute to
worker occupancy if execution capacity remains held.
}

\termrow{output write}{
The stage in which task output, state changes, generated data, or side effects
are written.
}

\termrow{commit}{
The stage in which task effects are finalized, persisted, checkpointed, or made
visible.
}

\termrow{acknowledge}{
The lifecycle point at which a worker confirms task completion to the queue or
broker.
}

\termrow{ack}{
A short form of acknowledge.
}

\termrow{discard}{
The lifecycle outcome in which a task instance or attempt is intentionally
dropped or ignored.
}

\termrow{dead-letter}{
The lifecycle outcome in which a task is moved to a dead-letter queue or
equivalent failure storage.
}

\termrow{finish}{
The lifecycle point at which a task attempt ends, whether by success, failure,
timeout, cancellation, or another terminal outcome.
}

\termrow{lease expiration}{
The lifecycle event in which a task reservation expires before successful
acknowledgement or completion.
}

\termrow{cancellation}{
The lifecycle event in which a task attempt is explicitly stopped before normal
completion.
}

\end{longtable}
